\scanpage{14}
\handwritten{\mspage{5}}
We assume \inserted{a) and b) satisfied for the induced morphism $i_{\bar s}:X_{\bar s}\handwritten{\longrightarrow}X_{\bar s}$} \deleted{that the induced morphism} for the fibers at $\bar s$ (\deleted{fixed} geometric point of $S$\deleted{, which it is in our interest to take} above the \uline{generic} point of $S$). \deleted{if $S$ is irreducible to have the maximum chance of \unsure{being able that|providing that} the serious hypothesis b). Let $z$ be a cycle} \unsure{a|\deleted{a}} \inserted{Let} $z\handwritten{\in}Z^n(Y)$ be such that the induced cycle on $Y_{\bar s}$ is algebraically equivalent to zero, and a fortiori cohomologically equivalent to zero. Thus \unsure{d|\deleted{l}d|\deleted{d}l}a class
\[
u(z)\handwritten{\in}H^1(S,R^{2n-1}f_{\handwritten{*}}(\underline{Z}_{\unsure{\chi|\ell}}(n))
\]
is defined. This being so, I say that \uline{this class is \deleted{a torsion class} contained in} $i^{\handwritten{*}}(H^{\deleted{\unsure{0}}1}(S,R^{2n-1}\deleted{\unsure{f|g}}\handwritten{g}_{\handwritten{*}}(\unsure{\underline{\deleted{Z}\handwritten{Q}}|\underline{Z}}_{\chi}(n)))$ \uline{modulo a torsion class}, at least if we \deleted{assume $S$ Japanese of dimension 1, or even} \deleted{\inserted{We apply 2) to $i_{\bar s}$}} \deleted{\inserted{if,}} allow ourselves to replace $S$ by a nonempty open subset(restriction \deleted{probably}\inserted{probably} unnecessary, but this statement will suffice for what we want to do with it).

To see this, we apply 2) to the morphism $i_{\bar s}$. The curve $C$ and the cycle $Z''$ can \unsure{d|\deleted{s}d|\deleted{d}s}e already be defined over a finite extension $K$ of $k(s)$, A trace argument allows us to assume that $K=k(s)$, and \deleted{if necessary taking the closure of the cycle} we can assume that $\unsure{\deleted{Z}C|\deleted{C}Z|\mathbb{C}}$ \deleted{and $Z$} is a proper and smooth relative curve over $S$ itself, and $Z''$ a cycle on $\unsure{\mathbb{X}|X|Y}\mathbin{\mathrm{x}}_S C$, \deleted{Then} and that the formula $z=i^{\handwritten{*}}(z')+z''$ holds on $\unsure{\deleted{X}Y|\deleted{Y}X|\mathbb{X}}\mathbin{\mathrm{x}}_S C$. It implies an analogous decomposition of $u(z)$, which is the sum of $i^{\handwritten{*}}(u(z'))$ and of $u(z'')$, so that it suffices to prove that $u(z'')$ is zero. Now $u(z'')$ induces by \deleted{d}construction 2) the zero homomorphism $H^1(C_{\bar s})\handwritten{\longrightarrow}H^{2n-1}(Y_{\bar s})$. I\unsure{m|\deleted{l}m|\deleted{m}l} therefore indeed follows from 1) that $u(z'')$ is zero, which was to be proved.

4.\quad Let $X$ be a projective and smooth scheme of dimension $m$\inserted{\handwritten{over algebraically closed $k$,}}, let us take a Lefschetz pencil, so as to obtain a blown-up variety no less projective and smooth $\overset{\handwritten{\sim}}{X}$ and a morphism
\[
\overset{\handwritten{\sim}}{f}:\overset{\handwritten{\sim}}{X}\handwritten{\longrightarrow}\deleted{S\mathrm{x}}P^1_k
\]
% The x in the preceding deletion records a visible extra cancellation cell, not a recovered old lexical x.
% The b in fibers has a long thin, partly interrupted extension of its ascender and an irregular retraced bowl. Pen retouch versus surviving type is uncertain; the mark is treated as part of this letter, not as an independent drawing or additional character.
whose fi\unsure{\handwritten{b}|b}res are the hyp\deleted{é}erplane sections of the pencil. Let $\unsure{B|S|U|\mathbb{S}}$ be a nonempty open subset of $P^1_k$ over which $\overset{\handwritten{\sim}}{X}$ is smooth, and let $\handwritten{\unsure{j|f}}:Y=\overset{\handwritten{\sim}}{X}\handwritten{|}B\handwritten{\longrightarrow\unsure{S}}$.
% archive sheet number: 13
